%% kapa: exam=1 round=34 year=2023 period=2 form=A subject=appraisal_law subject_ko=감정평가관계법규 seq=5 number=5
%% source: Q-Net 감정평가사 기출문제 (cst00309) artlSeq=5211789 file=raw/1008/5211789/2023년도 제34회 감정평가사 1차시험 2교시.pdf
%% conversion: kordoc 4.24.51 -> tools/gen1.py (auto); review=proofread-sonnet; flags=none
\begin{question}{5}
국토의 계획 및 이용에 관한 법령상 개발행위의 허가에 관한 설명으로 옳지 않은 것은? (단, 조례는 고려하지 않음)
\begin{choices}
  \item 개발행위허가를 제한하고자 하는 자가 시·도지사인 경우에는 당해 시·도에 설치된 지방도시계획위원회의 심의를 거쳐야 한다.
  \item 「사도법」에 의한 사도개설허가를 받은 토지의 분할은 개발행위허가를 받지 아니하고 할 수 있다.
  \item 시장 또는 군수는 개발행위허가에 조건을 붙이려는 때에는 미리 개발행위허가를 신청한 자의 의견을 들어야 한다.
  \item 「사방사업법」에 따른 사방사업을 위한 개발행위는 중앙도시계획위원회와 지방도시계획위원회의 심의를 거치지 아니한다.
  \item 개발행위허가를 받은 부지면적을 5퍼센트 확장하는 경우에는 별도의 변경허가를 받지 않아도 된다.
\end{choices}
\end{question}
\answer{5}
